\documentclass[11pt]{article}
\usepackage[a4paper,margin=28mm]{geometry}
\usepackage{amsmath,amssymb,amsthm,mathtools,hyperref}
\newcommand{\ddc}{dd^c}
\newcommand{\PP}{\mathbb P}
\newcommand{\MA}{\operatorname{MA}}
\newcommand{\I}{\mathcal I}
\newtheorem{theorem}{Theorem}
\newtheorem{lemma}[theorem]{Lemma}
\newtheorem{proposition}[theorem]{Proposition}
\theoremstyle{remark}\newtheorem{remark}[theorem]{Remark}
\title{Density-tail criteria for residual regularity\newline
in nonlogarithmic radial singularity types}
\author{}
\date{}
\begin{document}
\maketitle
\begin{abstract}
We give an explicit density criterion for regularity of the residual
after subtracting a prescribed nonlogarithmic singularity on projective
space. If the relative density error is $O(r^\delta)$, the radial
solution has a residual with exponent $\min\{1,2\delta\}$, extending
$C^1$ for $\delta\ge\frac12$ and $C^2$ for $\delta\ge1$, including
the critical values. Compact examples show that these derivative
thresholds cannot be lowered. Errors flat to every weighted derivative
order give smoothly flat residuals. The solutions have their actual
canonical poles and prescribed ordinary types, although all their
multiplier ideals are trivial. Nonnegative radial densities may vanish
away from the pole: positive cumulative central mass gives smooth
extension across those zeros. We construct compact examples with an
open zero annulus and a globally smooth residual. A conditional product construction
transfers the residual estimates to genuinely nonclosed backgrounds
with zero lower mass. No nonsymmetric uniqueness or general weighted
envelope regularity is asserted.
\end{abstract}

Use $d^c=\frac{i}{4\pi}(\bar\partial-\partial)$ and
$\ddc=\frac{i}{2\pi}\partial\bar\partial$.
All potentials have their genuine canonical representatives.
The ideal $\I(tu)$ uses local integrability of $|g|^2e^{-tu}$.
Ordinary equivalence means equality of singularities up to bounded
additive differences. Non-pluripolar products are taken by actual
bounded truncations with smooth-background corrections. The bounded
locality and multilinearity used below are those of
\cite[Section 2.1]{BGL}; we do not invoke its unrestricted product
theorem under a missing bounded mass hypothesis.

\section{The prescribed ordinary type}
Let $k\ge2$ and normalize the Fubini--Study form $\kappa$ on
$\PP^k$ by $\int\kappa^k=1$. In the chart $z_0=1$, set
\[
 r=\sum_{j=1}^k|z_j|^2,\qquad P=\frac r{1+r},\qquad
 T=1-\log P,
 \qquad u_\alpha=1-T^\alpha\quad(0<\alpha<1).
\]
This is a global function with its sole pole
$o=[1:0:\cdots:0]$; it is smooth elsewhere, including the hyperplane
$z_0=0$, where its value is zero. Put
\begin{equation}\label{eq:baseline}
 \Phi_\alpha=\log(1+r)+u_\alpha,
 \qquad Q_\alpha=r\Phi_\alpha'
      =P+\alpha(1-P)T^{\alpha-1}.
\end{equation}
Differentiation with respect to $P$ gives
\[
 \frac{dQ_\alpha}{dP}
 =1-\alpha T^{\alpha-1}
  +\alpha(1-\alpha)\frac{1-P}{P}T^{\alpha-2}>0.
\]
Thus $Q_\alpha>0$, $Q_\alpha'>0$ away from the pole.
The local radial Hessian has tangential eigenvalues
$\Phi_\alpha'=Q_\alpha/r$ and radial eigenvalue
$\Phi_\alpha'+r\Phi_\alpha''=Q_\alpha'$.
The decreasing radial truncations
$\max\{\Phi_\alpha(r),\Phi_\alpha(\varepsilon)\}$ are psh on a
small ball: their slopes as functions of $\log r$ are nonnegative
and increasing, with a nonnegative jump at $r=\varepsilon$.
Their limit is a genuine psh function, by
\cite[I.5.4]{Dem}, with canonical value $-\infty$ at the origin.
Hence $u_\alpha$ is globally $\kappa$-psh.

For any smooth radial local potential $\Phi$, with $Q=r\Phi'$,
radial integration of its determinant gives
\begin{equation}\label{eq:flux}
 \int_{a<|z|^2<b}(\ddc\Phi)^k=Q(b)^k-Q(a)^k.
\end{equation}
Indeed the radial integrand is proportional to
$kQ^{k-1}Q'\,dr$; the normalization constant equals one by
applying it to $\Phi=\log(1+r)$, whose total mass is one.
For the baseline, $Q_\alpha(0)=0$ and $Q_\alpha(\infty)=1$.
The actual non-pluripolar measure is therefore
$q_\alpha\kappa^k$ of total mass one, where
\begin{equation}\label{eq:density}
 q_\alpha(r)=\frac{(Q_\alpha^k)'(r)}{(P^k)'(r)}>0.
\end{equation}
Smoothness on the hyperplane and the exclusion of the canonical
pole in the product leave no further measure. Near $r=0$, with
$L=1-\log r$, direct differentiation gives
\begin{align}\label{eq:asymptotics}
 Q_\alpha&\sim\alpha L^{\alpha-1},&
 Q_\alpha'&\sim\frac{\alpha(1-\alpha)}rL^{\alpha-2},\\
 (Q_\alpha^k)'&\sim
 \frac{k\alpha^k(1-\alpha)}rL^{-1-k(1-\alpha)}.\nonumber
\end{align}
In particular $q_\alpha$ is in $L^1$ and in no $L^p$, $p>1$.

The types for different $\alpha$ are ordinarily distinct, since
their difference is unbounded at $o$. Yet every $\I(tu_\alpha)$
is trivial for finite $t>0$: integrability of the constant function
reduces near $o$ to
$\int^\infty e^{-kL+tL^\alpha}\,dL<\infty$.
Ordinary type information here is strictly finer than this ideal
information. We do not claim that $u_\alpha$ is an ordinary model.

\section{Density criterion and exact inverse}
Here ``radial'' means depending only on $r$ in the displayed chart.
Smooth radial data along the whole hyperplane have a smooth
expression in $t=1/r$ near $t=0$.

\begin{theorem}\label{thm:regularity}
Let $q\ge0$ be radial, smooth on $\PP^k\setminus\{o\}$, including
the hyperplane at infinity, and satisfy $\int q\kappa^k=1$.
Suppose that for some $\delta,A,r_0>0$,
\begin{equation}\label{eq:input}
 \left|\frac{q(r)}{q_\alpha(r)}-1\right|\le Ar^\delta
 \quad(0<r<r_0).
\end{equation}
Then there is a genuine radial $\kappa$-psh potential $U$, normalized
to have value zero on the hyperplane, with
\[
 \MA_\kappa(U)=q\kappa^k,
 \qquad U\sim u_\alpha.
\]
Its actual residual $R=U-u_\alpha$ off $o$ extends continuously
across $o$, with value $R_0$, and for small $r$ satisfies
\begin{equation}\label{eq:residual}
 \begin{aligned}
 |R-R_0|&\le Cr^\delta L^{\alpha-2},\\
 |R'|&\le Cr^{\delta-1}L^{\alpha-2},\qquad
 |R''|\le Cr^{\delta-2}L^{\alpha-2}.
 \end{aligned}
\end{equation}
It is locally $C^{0,\min\{1,2\delta\}}$.
For $\delta\ge\frac12$ it extends $C^1$ with zero gradient at $o$;
for $\delta\ge1$ it extends $C^2$ with zero Hessian at $o$.
No derivative hypothesis on the error in \eqref{eq:input} is needed
for these conclusions.
The potential and its residual are smooth on $\PP^k\setminus\{o\}$,
including every density zero there. The tail condition ensures positive
cumulative central mass; local density smoothness alone is not the criterion.
\end{theorem}

\begin{proof}
Let
\begin{equation}\label{eq:inverse}
 M(r)=\int_{|z|^2<r}q\kappa^k,
 \qquad Q(r)=M(r)^{1/k},
 \qquad U'(r)=\frac{Q(r)-P(r)}r,
 \quad U(\infty)=0.
\end{equation}
The tail condition makes $q$ positive near $o$, so $M(r)>0$ for every
$r>0$. One has $M'=q(P^k)'\ge0$, $Q'\ge0$ and $M(\infty)=1$.
Thus the positive root defining $Q$ is smooth for every finite $r>0$,
even where $q$ vanishes. This proves smoothness across all finite density
zeros without a vanishing-order assumption.
Smoothness of $q$ in $t=1/r$ implies that $M$ and $Q$ are smooth
in $t$ near zero, with $Q(0)=1$ in that coordinate. The equation
$dU/dt=-(Q-P)/t$ is smooth there, since $Q-P$ vanishes at
$t=0$. This both defines the normalization and proves smoothness
along the full hyperplane.

Write $b=k(1-\alpha)>0$. From \eqref{eq:input} and
\eqref{eq:asymptotics},
\begin{equation}\label{eq:tail}
 |M-Q_\alpha^k|
 \le C\int_0^r s^{\delta-1}(1-\log s)^{-1-b}\,ds
 \le Cr^\delta L^{-1-b}.
\end{equation}
The last inequality follows by writing $s=re^{-y}$ and bounding
$\int_0^\infty e^{-\delta y}(L+y)^{-1-b}\,dy$.
Dividing by $Q_\alpha^k\asymp L^{-b}$ shows that $Q/Q_\alpha$
tends to one; the mean value theorem for the positive $k$th root
then gives
\begin{equation}\label{eq:Qerror}
 |Q-Q_\alpha|\le Cr^\delta L^{\alpha-2}.
\end{equation}
Moreover the exact identity
\[
 \frac{Q'}{Q_\alpha'}=
 \frac q{q_\alpha}\left(\frac{Q_\alpha}{Q}\right)^{k-1}
\]
gives $|Q'-Q_\alpha'|\le Cr^{\delta-1}L^{\alpha-2}$.
Since $R'=(Q-Q_\alpha)/r$, integration to the origin gives a finite
$R_0$ and the first two inequalities in \eqref{eq:residual}.
Differentiating this formula gives the last inequality.

The local potential $\Phi=\log(1+r)+U$ has semipositive
Hessian away from $o$ because $Q/r>0$ and $Q'\ge0$. Also
$U=u_\alpha+R_0+o(1)$ near $o$, so $\Phi$ tends to $-\infty$.
The same radial truncation argument as above gives its genuine psh
extension, with the canonical pole retained. Thus $U$ is globally
$\kappa$-psh. Formula~\eqref{eq:flux} or the local determinant
ratio gives its measure $q\kappa^k$ away from $o$; the defining
non-pluripolar truncations exclude $o$ and are locally identical
to this smooth product elsewhere. The measure has exactly total
mass one. The bounded global residual proves the ordinary type.

In real coordinates, with $\rho=|z|=\sqrt r$,
\[
 |\nabla R|\le C\rho^{2\delta-1}L^{\alpha-2},
 \qquad |\nabla^2R|\le C\rho^{2\delta-2}L^{\alpha-2}.
\]
For $2\delta<1$, integration along radial radii bounds differences
by $C|\rho_1-\rho_2|^{2\delta}$ and hence by
$C|z_1-z_2|^{2\delta}$. For $2\delta\ge1$ the gradient is bounded,
giving local Lipschitz regularity. At $\delta=\frac12$ its bound
tends to zero and $R-R_0=o(\rho)$, which give the claimed $C^1$
extension. At $\delta=1$, the Hessian bound tends to zero and
$\nabla R=o(\rho)$, giving the $C^2$ extension. Larger $\delta$
only improves these bounds. These arguments justify crossing the
pole for $R$; they do not assign a finite value to $U(o)$.
\end{proof}

For completeness, among radial genuine solutions smooth off $o$
with this density, no central flux and no hyperplane singularity,
the potential is unique up to constants: its flux satisfies
$Q^k=M$, and \eqref{eq:inverse} then fixes its derivative.
We assert no uniqueness in a nonsymmetric family.

The precise relative stretching class is also verified. If
$a_* =\inf R$, the genuine $\kappa$-psh functions
$v_t=u_\alpha+t a_*$ satisfy, for every $t\ge1$,
\[
 t^{-1}v_t+(1-t^{-1})u_\alpha=u_\alpha+a_*\le U.
\]
At $t=1$ the term with zero coefficient is omitted. Hence
$U\in E_{\mathrm{env}}(\PP^k,\kappa,u_\alpha)$ in the actual
stretching definition, canonically at the common pole.
No envelope identity is substituted for that definition.

\section{Smooth residuals from flat errors}
\begin{proposition}\label{prop:flat}
Under the hypotheses of Theorem~\ref{thm:regularity}, suppose
$g=q/q_\alpha-1$ satisfies
\[
 (r\partial_r)^jg=O(r^N)
 \quad\text{for every }N,j\ge0.
\]
Then $R-R_0$ extends smoothly across $o$, with every derivative
zero there.
\end{proposition}
\begin{proof}
Put $x=-\log r$, $a(x)=Q_\alpha(e^{-x})^k$ and
$E(x)=M(e^{-x})-a(x)$. The exact central integral is
\[
 E(x)=\int_x^\infty g(e^{-y})[-a'(y)]\,dy,
 \qquad E'(x)=g(e^{-x})a'(x).
\]
Every derivative of $a$ has at most polynomial growth in $1+x$,
and $a\asymp(1+x)^{-b}$. The integral bounds $E$ by an
exponential of arbitrarily large rate. Repeated differentiation of
$E'$ expresses each derivative as a finite Leibniz sum of flat
derivatives of $g$ and polynomially bounded derivatives of $a$;
all derivatives of $E$ are therefore exponentially flat.
Every derivative of $1/a$ has at most polynomial growth, by
\eqref{eq:baseline}. Thus $E/a$ has the same flatness. Applying
smooth composition near one to
\[
 Q-Q_\alpha=Q_\alpha\bigl[(1+E/a)^{1/k}-1\bigr]
\]
proves flatness of every $x$ derivative of $Q-Q_\alpha$.
Since $dR/dx=-(Q-Q_\alpha)$, integration from $x$ to infinity
proves the same for $R-R_0$, including its zeroth derivative.
An ordinary $j$th derivative in $r$ is $r^{-j}$ times a polynomial
in $\partial_x$; arbitrarily strong exponential decay absorbs that
factor. All ordinary derivatives extend continuously as zero at
$r=0$, and repeated one-dimensional integration verifies smoothness
there. Composition with $r=|z|^2$ gives the conclusion.
\end{proof}

\section{Compact examples and sharp derivative thresholds}
Choose smooth nonnegative radial cutoffs $\chi_0,\chi_1$ with
disjoint supports: $\chi_0=1$ near $r=0$ and vanishes outside a
finite small ball; $\chi_1$ is nonzero in a finite annulus. Put
\begin{equation}\label{eq:examples}
 B=\frac{\int r^\delta\chi_0q_\alpha\kappa^k}
          {\int\chi_1q_\alpha\kappa^k},\qquad
 q_c=q_\alpha\bigl[1+c(r^\delta\chi_0-B\chi_1)\bigr].
\end{equation}
For sufficiently small $c>0$, $q_c$ is positive, smooth off $o$
and along the hyperplane, and has total mass exactly one. It is
comparable to $q_\alpha$ near $o$, hence is in no $L^p$, $p>1$.
Theorem~\ref{thm:regularity} gives an actual compact solution in
the prescribed ordinary type. Near $o$ its central flux difference is
exactly $c\int_0^r s^\delta(Q_\alpha^k)'\,ds$.
The substitution $s=re^{-y}$ and dominated convergence give
\[
 \int_0^r s^{\delta-1}(1-\log s)^{-1-b}\,ds
 \sim\delta^{-1}r^\delta L^{-1-b}.
\]
Using \eqref{eq:asymptotics} and the $k$th-root expansion yields
\begin{align*}
 R_c'(r)&\sim\frac{c\alpha(1-\alpha)}\delta
                  r^{\delta-1}L^{\alpha-2},\\
 R_c(r)-R_c(0)&\sim\frac{c\alpha(1-\alpha)}{\delta^2}
                  r^\delta L^{\alpha-2}.
\end{align*}
For $\delta<\frac12$, the radial gradient
$2\sqrt r\,R_c'$ diverges, ruling out $C^1$ extension. For
$\delta<1$, a tangential real Hessian eigenvalue $2R_c'$ diverges,
ruling out $C^2$ extension. The annular compensation changes no
central flux. For $2\delta<1$ the displayed asymptotic also rules
out any larger H\"older exponent. The critical derivative thresholds
in Theorem~\ref{thm:regularity} are therefore sharp for this family.

Replacing $r^\delta\chi_0$ in \eqref{eq:examples} by
$e^{-1/r}\chi_0$ and adjusting $B$ by the same integral formula
gives a nonzero mass-preserving perturbation satisfying every
flat-error hypothesis. Proposition~\ref{prop:flat} then supplies
an actual smoothly flat residual across the pole.

\section{Smooth residuals with an open density-zero annulus}
\begin{proposition}\label{prop:zeros}
For every $k\ge2$ and $0<\alpha<1$, there is a nonnegative radial
probability density $q_0$, smooth off $o$ and along the hyperplane,
which equals $q_\alpha$ near $o$ and vanishes on a nonempty open
annulus. Its actual solution has ordinary type $u_\alpha$ and a
globally smooth residual, constant near $o$.
\end{proposition}
\begin{proof}
Choose a smooth radial $h$ with $0\le h\le1$, equal to zero on
a nonempty finite annulus and equal to one near $o$ and outside
a larger finite annulus. Choose a nonzero nonnegative smooth radial
cutoff $\chi$ supported in a disjoint finite annulus where $h=1$.
Set
\[
 C=\frac{\int(1-h)q_\alpha\kappa^k}{\int\chi q_\alpha\kappa^k}>0,
 \qquad q_0=q_\alpha(h+C\chi).
\]
This is nonnegative and has integral one. It is smooth at every
point off $o$, including the hyperplane, and equals $q_\alpha$
near $o$. It vanishes on the chosen open annulus and is still in
no $L^p$, $p>1$. Theorem~\ref{thm:regularity} applies, with
zero relative density error on a neighborhood of $o$. There the
exact central fluxes agree, so $Q=Q_\alpha$ and $R'=0$.
Hence $R$ is constant near $o$ and extends smoothly across it.
It is smooth everywhere else by the positive-root inverse, including
the entire zero annulus. The solution retains its genuine pole;
this smoothness assertion concerns the residual.
\end{proof}

\section{Conditional nonclosed products and limits of scope}
Let $(X_0,h_0)$ be any connected compact Hermitian $m$-fold with
$v_-(h_0)=0$, and write $V_0=\int_{X_0}h_0^m>0$.
On $X=X_0\times\PP^k$ put $N=m+k$ and
$H=\operatorname{pr}_0^*h_0+\operatorname{pr}_Y^*\kappa$.
The actual pullback of $U$ obeys
\[
 \MA_H(U)=\binom Nm\operatorname{pr}_0^*(h_0^m)
                 \wedge\operatorname{pr}_Y^*(q\kappa^k),
 \qquad \int_X\MA_H(U)=\binom Nm V_0.
\]
Indeed the bounded truncation $\max\{U,-j\}$ is a pullback.
The bounded corrected-product expansion leaves only the displayed
binomial term by dimensions; multiplication by $1_{\{U>-j\}}$
and increasing $j$ give the exact formula and finite mass.
The residual is the actual pullback of $R$, so its proved regularity
extends across $X_0\times\{o\}$. It is ordinarily equivalent to the
pullback of $u_\alpha$ and belongs to the corresponding relative
stretching class by the same explicit $v_t$.

Also $v_-(H)=0$: bounded $h_0$-psh potentials with masses tending
to zero have pullbacks of mass $\binom Nm$ times those masses.
The background is genuinely nonclosed, since a closed positive
$h_0$ would have all bounded masses equal to $V_0>0$ by Stokes.
This conclusion retains the factor hypothesis as an antecedent;
no named-factor theorem or exact individual normalization threshold
is imported here.

These are radial density-tail conditions in specific genuinely
prescribed nonlogarithmic ordinary types. They neither establish
regularity in arbitrary nonclosed types nor settle general density
zeros or weighted extremals. The proved radial geometry permits smooth
crossing of its density zeros and quantitative crossing of the pole
for the residual. The full S06 question remains open.

\begin{thebibliography}{9}
\bibitem{Dem} J.-P. Demailly,
\emph{Complex Analytic and Differential Geometry}, 21 June 2012,
I.5.4 and I.5.8, pp. 39--40;
\url{https://people.math.harvard.edu/~demarco/Math274/Demailly_ComplexAnalyticDiffGeom.pdf}.
\bibitem{BGL} S. Boucksom, V. Guedj and C. H. Lu,
\emph{Volumes of Bott--Chern classes},
\href{https://arxiv.org/abs/2406.01090}{arXiv:2406.01090},
Section 2.1, Proposition 2.1 and its complete proof, pp. 12--14.
\end{thebibliography}
\end{document}
